\documentclass[11pt]{article}
\usepackage{amsmath,amssymb,amsthm}
\usepackage[margin=1in]{geometry}
\usepackage{algorithm}
\usepackage{algpseudocode}
\usepackage{xcolor}

\newtheorem{theorem}{Theorem}
\newtheorem{lemma}{Lemma}
\newtheorem{claim}{Claim}
\newtheorem{remark}{Remark}
\newtheorem{assumption}{Assumption}

\DeclareMathOperator{\Null}{Null}
\DeclareMathOperator{\Range}{Range}

\title{The Recursive Complete Minimum-Norm HSS Solver}
\author{Arjun Sethi-Olowin}
\date{\today}

\begin{document}
\maketitle

\section{Introduction}

These notes provide the notation and the algorithm to understand a direct
hierarchically semi-separable (HSS) ULV-based minimum norm solver.

\section{Setup and Notation}

We are considering $H\in\mathbb{C}^{m\times n}$ as a matrix that admits HSS
structure, with $m<n$. We wish to solve
\[
x^* = \operatorname*{arg\,min}_{Hx=b}\|x\|_2 = H^*(HH^*)^{-1}b.
\]
At a generic sibling pair $\tau,\sigma\in\{1,2\}$ ($\tau\neq\sigma$), the diagonal
and off-diagonal blocks are
\[
D_\tau\in\mathbb{C}^{n_\tau\times p_\tau},\quad p_\tau>n_\tau,\qquad
U_\tau\in\mathbb{C}^{n_\tau\times k_\tau},\qquad
V_\tau\in\mathbb{C}^{\ell_\tau\times p_\tau},
\]
\[
H_{\tau\sigma} = U_\tau B_{\tau\sigma}V_\sigma^*,\qquad B_{\tau\sigma}\in\mathbb{C}^{k_\tau\times\ell_\sigma},
\]
with $k_\tau$ the rank of $\tau$'s own row generator and $\ell_\tau$ the rank of
$\tau$'s own column generator.

For the algorithm of Section~\ref{sec:algorithm} we require:
\begin{assumption}
\label{as:slack}
Every leaf $\tau$ has enough slack for its sibling's coupling rank:
$\ell_\tau + n_\tau \le p_\tau$.
\end{assumption}
This is a sizing condition, not a compression-quality one: the Size Reduction
step (Section~\ref{subsec:sizereduction}) LQ-factors the $p_\tau'\times p_\tau$
stack $[V_\tau;D_\tau]$, $p_\tau':=\ell_\tau+n_\tau$, and an LQ factorization can
only produce a full $p_\tau'\times p_\tau'$ unitary factor -- as opposed to a
truncated one that would silently lose information -- when $p_\tau'\le p_\tau$.
Assumption~\ref{as:slack} is exactly this condition.

\section{The Recursive Minimum-Norm Solver}
\label{sec:algorithm}

At a leaf-sibling pair $\tau,\sigma$, two transformations reduce $Hx=b$ to a
smaller system of the same form: Size Reduction, which uses the joint
structure of $V_\tau$ and $D_\tau$ to shrink the column count of each leaf, and
Decoupling, which uses two further unitary transforms per leaf to split off a
forced, uniquely-determined block of $x$ from a remaining, still-coupled
block. Once every leaf pair in the tree has been processed this way, what
remains is again exactly an HSS system, one level shallower, on which the same
procedure recurses; the recursion terminates once the tree has collapsed to a
single dense leaf, solved directly by the usual normal-equations formula.

\subsection{Size Reduction}
\label{subsec:sizereduction}

For each leaf $\tau$, LQ-factor the stacked matrix
\[
\begin{bmatrix}V_\tau\\ D_\tau\end{bmatrix}
=
\begin{bmatrix}\tilde V_\tau & 0\\ \tilde D_\tau & 0\end{bmatrix}\Omega_\tau^*,
\qquad
\tilde V_\tau\in\mathbb{C}^{\ell_\tau\times p_\tau'},\quad
\tilde D_\tau\in\mathbb{C}^{n_\tau\times p_\tau'},\quad
p_\tau' := \ell_\tau+n_\tau,
\]
with $\Omega_\tau\in\mathbb{C}^{p_\tau\times p_\tau}$ unitary. $V_\tau$ and $D_\tau$
share the same discarded right null-space directions by construction, since
they are LQ-factored \emph{jointly} -- this is what makes
Remark~\ref{rem:sizereduction} below go through.

Discarding the zero columns gives the size-reduced pair
\[
\tilde H =
\begin{bmatrix}
\tilde D_1 & U_1B_{12}\tilde V_2^*\\
U_2B_{21}\tilde V_1^* & \tilde D_2
\end{bmatrix}.
\]

\begin{remark}
\label{rem:sizereduction}
If $\tilde x$ solves $\tilde Hx=b$ with minimum norm, then
$x^*=\Omega\begin{bmatrix}\tilde x\\0\end{bmatrix}$, $\Omega=\operatorname{diag}(\Omega_1,\Omega_2)$,
is the minimum-norm solution of $Hx=b$. (Proved as Lemma~\ref{lem:orthogonal} below.)
\end{remark}

\subsection{Decoupling}
\label{subsec:decoupling}

\paragraph{Step 1: QL of $U_\tau$.}
\[
U_\tau = P_\tau\begin{bmatrix}0\\ \hat U_\tau\end{bmatrix},\qquad
P_\tau\in\mathbb{C}^{n_\tau\times n_\tau}\text{ unitary},\quad
\hat U_\tau\in\mathbb{C}^{k_\tau\times k_\tau}\text{ lower triangular}.
\]
Apply $P_\tau^*$ on the left to $\tilde D_\tau$ (and to $\tau$'s rows of $b$):
$\tilde D_\tau\leftarrow P_\tau^*\tilde D_\tau$. This zeros the top $t_\tau:=n_\tau-k_\tau$
rows of $U_\tau$, densifying $\tilde D_\tau$.

\paragraph{Step 2: LQ of the top rows of $P_\tau^*\tilde D_\tau$.}
Partition $P_\tau^*\tilde D_\tau=\begin{bmatrix}\tilde D_{\tau;11} & \tilde D_{\tau;12}\\ \tilde D_{\tau;21}&\tilde D_{\tau;22}\end{bmatrix}$,
top block $t_\tau$ rows. LQ-factor the top block row:
\[
\begin{bmatrix}\tilde D_{\tau;11}&\tilde D_{\tau;12}\end{bmatrix}
=\begin{bmatrix}\hat D_{\tau;11}&0\end{bmatrix}Q_\tau^*,
\qquad
Q_\tau\in\mathbb{C}^{p_\tau'\times p_\tau'}\text{ unitary},\quad
\hat D_{\tau;11}\in\mathbb{C}^{t_\tau\times t_\tau}\text{ lower triangular}.
\]
Apply $Q_\tau^*$ on the right to the rest of $\tilde D_\tau$ and to $\tau$'s own
$\tilde V_\tau^*$ (they occupy the same reduced column space):
$\tilde D_\tau\leftarrow \tilde D_\tau Q_\tau^*$, $\tilde V_\tau^*\leftarrow \tilde V_\tau^*Q_\tau^*$.

\paragraph{Result.}
After forming $\hat H=P^*\tilde HQ^*$, $\hat b=P^*b$, the top $t_\tau$ rows of
leaf $\tau$'s block row read $\begin{bmatrix}\hat D_{\tau;11}&0&0&0\end{bmatrix}$
-- decoupled from the other leaf's columns. The bottom $k_\tau$ rows (over the
remaining $\ell_\tau+k_\tau$ unknowns) remain dense and coupled.

\subsection{Forced Solve and the Reduced System}
\label{subsec:forcedsolve}

Write $\hat x_\tau=\begin{bmatrix}X_\tau^{(1)}\\X_\tau^{(2)}\end{bmatrix}$,
$X_\tau^{(1)}\in\mathbb{C}^{t_\tau}$ decoupled, $X_\tau^{(2)}\in\mathbb{C}^{\ell_\tau+k_\tau}$ still coupled.
The decoupled top rows give
\[
\hat D_{\tau;11}X_\tau^{(1)} = \hat b_{\tau;1},
\]
solved directly by forward/back substitution -- unique, no freedom.

Substituting $X_\tau^{(1)}$ into the bottom rows gives the reduced system
\[
\underbrace{\begin{bmatrix}
\hat D_{1;22} & \hat U_1B_{12}\hat V_2^*\\
\hat U_2B_{21}\hat V_1^* & \hat D_{2;22}
\end{bmatrix}}_{=:\bar H}
\begin{bmatrix}X_1^{(2)}\\X_2^{(2)}\end{bmatrix}
=
\underbrace{\begin{bmatrix}\hat b_{1;2}-\hat D_{1;21}X_1^{(1)}\\ \hat b_{2;2}-\hat D_{2;21}X_2^{(1)}\end{bmatrix}}_{=:\bar b}.
\]

\begin{remark}
$\bar H$ is again exactly HSS: $\hat D_{\tau;22}$ is (in general) still wide,
and the off-diagonal blocks retain the $U\cdot B\cdot V^*$ form untouched by
this step's transforms. Across a full tree level, all sibling pairs reduce
simultaneously to blocks that assemble into a single HSS matrix with one less
level.
\end{remark}

\subsection{Reconstruction}
\label{subsec:reconstruction}

Given $\bar x=\begin{bmatrix}X_1^{(2)}\\X_2^{(2)}\end{bmatrix}$ from the
recursive call,
\[
\hat x_\tau = \begin{bmatrix}X_\tau^{(1)}\\X_\tau^{(2)}\end{bmatrix},\qquad
\tilde x_\tau = Q_\tau^*\hat x_\tau,\qquad
x_\tau = \Omega_\tau\begin{bmatrix}\tilde x_\tau\\0\end{bmatrix}.
\]

\begin{remark}
Depending on which LQ/QL convention an implementation's helper routines use
($A=LQ$ vs.\ $A=QL$, and which side $Q$ is applied from), the collected factors
may need to be applied transposed relative to the equations above -- check
against whatever convention is actually in force before trusting the signs.
\end{remark}

\section{Proof of Optimality}

\subsection{Two invariance lemmas}

\begin{lemma}[Right unitary substitution]
\label{lem:orthogonal}
Let $\Omega$ be unitary and suppose $A\Omega=[\tilde A,0]$ (i.e.\ the last
columns of $A\Omega$ vanish). If $\tilde x^*$ is the minimum-norm solution of
$\tilde A\tilde x=b$, then $x^*=\Omega\begin{bmatrix}\tilde x^*\\0\end{bmatrix}$ is
the minimum-norm solution of $Ax=b$.
\end{lemma}
\begin{proof}
Since $A\Omega=[\tilde A,0]$, for any $\tilde x,w$,
\[
A\Omega\begin{bmatrix}\tilde x\\w\end{bmatrix} = [\tilde A,0]\begin{bmatrix}\tilde x\\w\end{bmatrix} = \tilde A\tilde x,
\]
independent of $w$. So if $\tilde A\tilde x=b$, then $x:=\Omega\begin{bmatrix}\tilde x\\w\end{bmatrix}$
satisfies $Ax=b$ for \emph{every} choice of $w$. Conversely, since $\Omega$ is
unitary (hence invertible), every $x$ can be written this way: set
$\begin{bmatrix}\tilde x\\w\end{bmatrix}:=\Omega^*x$. So the feasible set of
$Ax=b$ is exactly
\[
\left\{\Omega\begin{bmatrix}\tilde x\\w\end{bmatrix} \;:\; \tilde A\tilde x=b,\ w\text{ arbitrary}\right\}.
\]

Because $\Omega$ is unitary, $\|x\|_2=\left\|\Omega\begin{bmatrix}\tilde x\\w\end{bmatrix}\right\|_2
=\sqrt{\|\tilde x\|_2^2+\|w\|_2^2}$. For any fixed feasible $\tilde x$, this is
minimized over the free variable $w$ by $w=0$, giving $\|x\|_2=\|\tilde x\|_2$.
Minimizing what remains over feasible $\tilde x$ (i.e.\ over $\tilde A\tilde x=b$)
is exactly the definition of $\tilde x^*$. Hence the minimum-norm $x$ is
$x^*=\Omega\begin{bmatrix}\tilde x^*\\0\end{bmatrix}$.
\end{proof}

\begin{lemma}[Two-sided unitary substitution]
\label{lem:twosided}
Let $P,Q$ be unitary and $\hat A=P^*AQ^*$, $\hat b=P^*b$. If $\hat x^*$ is the
minimum-norm solution of $\hat A\hat x=\hat b$, then $x^*=Q^*\hat x^*$ is the
minimum-norm solution of $Ax=b$.
\end{lemma}
\begin{proof}
Since $\hat A=P^*AQ^*$ and $P$ is unitary, $AQ^*=P\hat A$. So for any $\hat x$,
setting $x:=Q^*\hat x$ gives
\[
Ax = AQ^*\hat x = P\hat A\hat x.
\]
If $\hat A\hat x=\hat b$, then $Ax=P\hat b=PP^*b=b$ (using $\hat b=P^*b$ and $P$
unitary), so $x:=Q^*\hat x$ is feasible for $Ax=b$. Conversely, since $Q$ is
unitary (hence invertible), every $x$ can be written this way: set
$\hat x:=Qx$, and then $Ax=b\iff P^*Ax=P^*b\iff P^*AQ^*\hat x=\hat b\iff \hat A\hat x=\hat b$.
So $\hat x\mapsto x=Q^*\hat x$ is a bijection between the feasible sets of
$\hat A\hat x=\hat b$ and $Ax=b$.

Because $Q$ is unitary, $\|x\|_2=\|Q^*\hat x\|_2=\|\hat x\|_2$ for every $\hat x$.
So minimizing $\|x\|_2$ over $Ax=b$ is the same problem, under this
norm-preserving bijection, as minimizing $\|\hat x\|_2$ over $\hat A\hat x=\hat b$,
and the minimizers correspond: $x^*=Q^*\hat x^*$.
\end{proof}

\subsection{One-level case}

\begin{theorem}
\label{thm:onelevel}
Under Assumption~\ref{as:slack} and $\hat D_{\tau;11}$ nonsingular for $\tau=1,2$,
the vector $x$ produced by the construction of Section~\ref{sec:algorithm} with
$\bar x$ taken as the minimum-norm solution of $\bar H\bar x=\bar b$ is the
minimum-norm solution of $Hx=b$.
\end{theorem}
\begin{proof}
By Lemma~\ref{lem:orthogonal} (Size Reduction, Section~\ref{subsec:sizereduction},
applied once per leaf, jointly on $V_\tau,D_\tau$), the minimum-norm solution of
$Hx=b$ is $\Omega[\tilde x^*;0]$ where $\tilde x^*$ minimizes $\|\tilde x\|_2$
over $\tilde H\tilde x=b$. By Lemma~\ref{lem:twosided} (Decoupling,
Section~\ref{subsec:decoupling}), the minimum-norm solution of
$\tilde H\tilde x=b$ is $Q^*\hat x^*$ where $\hat x^*$ minimizes $\|\hat x\|_2$
over $\hat H\hat x=\hat b$.

It remains to identify $\hat x^*$. Every feasible $\hat x$ (not just the
minimum-norm one) satisfies $\hat D_{\tau;11}X_\tau^{(1)}=\hat b_{\tau;1}$
exactly, and $\hat D_{\tau;11}$ is square and nonsingular, so $X_\tau^{(1)}$ is
the \emph{same} for every feasible point. Hence
\[
\|\hat x\|_2^2 = \|X_1^{(1)}\|_2^2+\|X_2^{(1)}\|_2^2+\|X_1^{(2)}\|_2^2+\|X_2^{(2)}\|_2^2
\]
is minimized over the feasible set exactly by minimizing the free part
$\|X_1^{(2)}\|_2^2+\|X_2^{(2)}\|_2^2$ subject to $\bar H\bar x=\bar b$ (the
constraint remaining once $X_\tau^{(1)}$ is fixed and moved to the right-hand
side) -- i.e.\ $\bar x^*=(X_1^{(2)*},X_2^{(2)*})$ is exactly the minimum-norm
solution of $\bar H\bar x=\bar b$. Composing the three substitutions gives the
claimed $x$.
\end{proof}

\subsection{General recursive case}

\begin{theorem}
\label{thm:general}
For an $L$-level HSS matrix satisfying Assumption~\ref{as:slack} at every leaf
pair encountered, with every $\hat D_{\tau;11}$ nonsingular, the recursive
algorithm of Section~\ref{sec:algorithm} returns the minimum-norm solution of
$Hx=b$.
\end{theorem}
\begin{proof}
Induction on $L$. Base case $L=0$ ($H$ a single dense leaf): the algorithm
returns $H^*(HH^*)^{-1}b$, which is exactly the minimum-norm solution.

Inductive step: assume the recursive call on the $(L-1)$-level reduced matrix
$\bar H$ returns its minimum-norm solution $\bar x^*$. The argument of
Theorem~\ref{thm:onelevel} did not use $L=1$ anywhere except in asserting that
$\bar x^*$ is the minimum-norm solution of $\bar H\bar x=\bar b$; replacing
that assertion with the inductive hypothesis, the same proof goes through
verbatim and shows the reconstructed $x$ is the minimum-norm solution of
$Hx=b$.
\end{proof}

\begin{remark}
The two-level case is the first instance where this induction is doing real
work: the one-level reduced system is small enough to solve directly, but at
two levels $\bar H$ is itself a genuine one-level HSS matrix, and
Theorem~\ref{thm:onelevel} applied to it is what Theorem~\ref{thm:general}
needs at the base of the induction proper.
\end{remark}

\section{Open Items}

\begin{itemize}
\item Assumption~\ref{as:slack} and the nonsingularity of $\hat D_{\tau;11}$ --
  when do these actually hold for an HSS matrix arising from a genuine
  approximation (as opposed to being assumed outright)?
\item Complexity analysis: flop count per level, and the total over $L$ levels.
\item What happens when Assumption~\ref{as:slack} fails at some leaf pair --
  does the proof need anything beyond ``it's now a shallower instance of the
  same problem'' once that pair is densified away?
\end{itemize}

% Algorithm pseudocode, kept separate from the step-by-step write-up above so
% it can be cut entirely if Section~\ref{sec:algorithm} reads clearly without it.
\section{Algorithm Summary}

\begin{algorithm}[h]
\caption{Recursive Complete Minimum-Norm HSS Solve}
\label{alg:minnorm}
\begin{algorithmic}[1]
\Require HSS matrix $H\in\mathbb{C}^{m\times n}$, $m\le n$, full row rank; $b\in\mathbb{C}^m$
\If{$H$ is a single dense leaf}
  \State \Return $x = H^*(HH^*)^{-1}b$ \Comment{via \texttt{pinv}/\texttt{backslash}}
\Else
  \State At every leaf-sibling pair $(\tau,\sigma)$ in the tree, simultaneously:
  \State \quad Size Reduction: $\Omega_\tau$ (\S\ref{subsec:sizereduction})
  \State \quad Decoupling: $P_\tau,\,Q_\tau$ (\S\ref{subsec:decoupling})
  \State \quad Forced solve for $X_\tau^{(1)}$ (\S\ref{subsec:forcedsolve})
  \State Assemble the remaining free blocks $\{X_\tau^{(2)}\}$ into a reduced HSS
         matrix $\bar H$ one level shallower, with reduced right-hand side $\bar b$
  \State $\bar x \gets$ \textsc{RecursiveMinNormSolve}$(\bar H,\bar b)$ \Comment{recurse}
  \State Reconstruct $x$ from $\bar x$ and $\{\Omega_\tau,P_\tau,Q_\tau\}$ (\S\ref{subsec:reconstruction})
  \State \Return $x$
\EndIf
\end{algorithmic}
\end{algorithm}

\end{document}
